\documentclass[12pt,a4paper,oneside]{article}

\usepackage[spanish]{babel}

\def\mydate{\leavevmode\hbox{\the\year-\twodigits\month-\twodigits\day}}
\def\twodigits#1{\ifnum#1<10 0\fi\the#1}

\usepackage{fancyhdr}
\usepackage{geometry}
\usepackage{graphicx}
\usepackage{wrapfig}
\usepackage{lastpage}
\usepackage[hidelinks]{hyperref}
\usepackage{authblk}
\usepackage{bookmark}

\usepackage[utf8]{inputenc} % Required for inputting international characters
\usepackage[T1]{fontenc} % Output font encoding for international characters
\usepackage{mathpazo} % Palatino font

\makeatletter
\newcommand{\subtitledoc}[1]{\newcommand{\@subtitledoc}{#1}}
\newcommand{\instituto}[1]{\newcommand{\@instituto}{#1}}
\newcommand{\carrera}[1]{\newcommand{\@carrera}{#1}}
\newcommand{\professor}[1]{\newcommand{\@professor}{#1}}
\newcommand{\catedra}[1]{\newcommand{\@catedra}{#1}}
\newcommand{\curso}[1]{\newcommand{\@curso}{#1}}
\newcommand{\legajo}[1]{\newcommand{\@legajo}{#1}}
\newcommand{\footerauthor}[1]{\newcommand{\@footerauthor}{#1}}
\newcommand{\footerlegajo}[1]{\newcommand{\@footerlegajo}{#1}}
\newcommand{\footercatedra}[1]{\newcommand{\@footercatedra}{#1}}

%Configuracion de hoja (margenes y tamaño)
\geometry{a4paper,margin=1in}
\setlength\headheight{28pt}

%Metadata para pdf
%\hypersetup{
%    pdftitle={\@title},
%    pdfsubject={\@catedra\ - \@subtitledoc},
%    pdfauthor={\@author\ - \@legajo\ - \@curso},
%    pdfkeywords={\@title\ \@catedra\ \@subtitledoc\ \@author\ \@legajo\ \@curso}
%}

%formato de encabezado y pie para todas las paginas.
\fancyhead[L]{
  \begin{minipage}[b]{7.5mm}
    \includegraphics[width=7mm]{imagenes/COPYlogo-utn.eps}
  \end{minipage}
  \begin{minipage}[b]{85mm}
    \textbf{Alumnos: }\@footerauthor\\
    \textbf{Legajos: }\@footerlegajo
  \end{minipage}
}
\fancyhead[R]{
  %\textbf{Comisión:} \@curso 
  \textbf{Fecha:} \mydate \\
  \textbf{Cátedra:} \@footercatedra \text{ (}\@curso\text{)}
}
\fancyfoot[C]{} %eliminar antiguo numero de pagina
\fancyfoot[R]{Página \thepage\ de \pageref{LastPage}}
\renewcommand{\headrulewidth}{0.5pt}
\renewcommand{\footrulewidth}{0.5pt}
\pagestyle{fancy}

\addto\captionsspanish{%
  \renewcommand{\contentsname}%
  {CONTENIDO}%
}

\renewcommand{\maketitle}{%
  \newpage
  \thispagestyle{empty}

  \begin{center}

    \textsc{\LARGE \@instituto}\par
    \vspace{0.5cm}
    \textsc{\Large \@carrera}\par
    \vspace{1cm}
    \includegraphics[width=0.30\textwidth]{imagenes/COPYlogo-utn.eps}\par
    \vspace{1cm}

    \textsc{\Large \@catedra}\par
    \vspace{1.5cm}

    \hrule
    \vspace{0.5cm}
    {\huge\bfseries \@title}\par
    \vspace{0.4cm}
    \textsc{\large \@subtitledoc}
    \vspace{0.5cm}
    \hrule \vspace{\baselineskip}


  \end{center}

  \vspace{2cm}

  {\noindent
    \begin{minipage}[t]{.2\textwidth}
      \raggedright
      \textbf{DOCENTES} \par
      \vspace{\baselineskip}
      %\textbf{CÁTEDRA} \par
      \medskip
      \medskip
      \textbf{COMISIÓN} \par
      \medskip
      \medskip
      \textbf{ALUMNOS} \par
    \end{minipage}%
    \begin{minipage}[t]{.05\textwidth}
      \raggedright
      %\textbf{:} \par
      %~\\
      %~\\
      %\textbf{:} \par
      %\textbf{:} \par
      %\medskip
      %\medskip
      %\textbf{:} \par
    \end{minipage}%
    \begin{minipage}[t]{.55\textwidth}
      \raggedright
      \@professor. \par
      \medskip
      %\@catedra. \par
      \medskip
      \@curso \par
      \medskip
      \medskip
      \@author \par
    \end{minipage}%
    \begin{minipage}[t]{.15\textwidth}
      \raggedright
      \vspace{\baselineskip}
      \vspace{\baselineskip}
      \medskip
      \bigskip
      % \medskip
      \medskip
      \medskip
      \@legajo \par
    \end{minipage}
  }
  \vfill
  \begin{center}
    {\Large\textbf{Córdoba, \@date}} \vspace{\baselineskip}
    %\medskip
    %\textbf{Córdoba, Argentina} \par
    \newpage
  \end{center}
}

\makeatother

\usepackage[nottoc,numbib]{tocbibind}
\usepackage{amsmath}
\usepackage{siunitx}
\usepackage{amssymb}
\usepackage{xfrac}
\usepackage{multirow}
\usepackage{multicol}
\usepackage{float}
\usepackage{circuitikz}
\usepackage{tikz}
\usepackage{pgfplots}
\usepgfplotslibrary{groupplots}
\usepackage{float}

\pgfuseplotmark{square*}
\pgfuseplotmark{triangle*}
\pgfplotsset{compat=1.18}
\usetikzlibrary{calc, positioning}


\instituto{Universidad Tecnológica Nacional\\Facultad Regional Córdoba}
\carrera{Ingeniería Electrónica}
\title{Titulo}
\subtitledoc{Subtitulo}
\catedra{Catedra}
\professor{XXXXXXXXXX XXXXXXXX. \par XXXXXXXXXX XXXXXXXX.}
\curso{XRX}
\author{XXXXX XXXXX, XXXXX XXXXX. \par XXXXX XXXXX, XXXXX XXXXX.}
\legajo{XXXXX\par XXXXX}
\footerauthor{XXXXX, XXXXX}
\footerlegajo{XXXXX, XXXXX}
\footercatedra{XXXX}
\date{\today}

\newcommand\plotScale{0.7}

\begin{document} %Maximo 10 paginas
\maketitle
\tableofcontents
\newpage
\section{Introducción}
\section{Marco teorico}
\section{Primera Parte}
\subsection{Circuito}
\begin{center}
\begin{circuitikz}[american]
  \draw (0,0) to[V, l=$V_{CC}$, invert] ++(0,4)
  to[short] ++(3,0) 
  to[R=1\unit{\kilo\ohm}/2W] ++(0,-2)
  to[biD*,a=DB3, v^=$V_{A_1A_2}$] ++(0,-2)
  to[short] (0,0) node[ground]{};
\end{circuitikz}
\end{center}
\subsection{Procedimiento}
\begin{enumerate}
  \item Armar el circuito seleccionando un correcto valor de R en función del
    datasheet del DIAC.
  \item Variar la tensión de alimentación (V1) desde 0V a 50V según la tabla que se
    observa aquí abajo.
  \item Medir la corriente y caída de tensión en el DIAC.
  \item Invertir los terminales del DIAC y repetir las variaciones y mediciones
    expresadas en el punto 1 y 2
\end{enumerate}
\subsection{Calculo de R} %iP = 0.3A (50-14)/0.3  (ST datasheet) DB3
\subsection{Simulación}
\begin{figure}[H]
  \begin{center}
    %\includegraphics[width=0.5\textwidth, clip, trim=0.5cm 10cm 0.5cm 8cm]{inc/sim/circuits/CircuitoSegundaParte.pdf}
  \end{center}
  \caption{Circuito simulado}
  \label{fig:CircuitoPrimeraParte}
\end{figure}
\begin{center}
  \begin{tikzpicture}[
    ]
    \begin{axis}[
      width=\plotScale\linewidth,
      xlabel=$V_{A_1A_2}$, 
      ylabel=$I$, 
      grid=both,
      thick
      ]
      \addplot[color=red,
        skip coords between index={0}{71},
        ] table [x index=1, y index=3]{./inc/sim/PrimeraParte.txt};
      \addplot[color=blue,
        skip coords between index={69}{142},
        ] table [x index=1, y index=3]{./inc/sim/PrimeraParte.txt};
    \end{axis}
  \end{tikzpicture}
\end{center}
\subsection{Experimental}
% \begin{table}[H]
%   \begin{center}
%     \begin{tabular}{c|c|c}
%       $V_{CC}$ &$V_{A_1A_2}$ &$I$  \\
%       \hline
%       5     &5    &0     \\
%       10    &10   &1   uA   \\
%       15    &15   &1.5 uA \\
%       20    &20   &2   uA   \\
%       22    &22   &2.2 uA \\
%       24    &24   &2.4 uA \\
%       26    &26   &2.6 uA \\
%       28    &28   &2.7 uA \\
%       30    &30   &3   uA   \\
%       32    &32   &3.2 uA \\
%       33.4  &27.2 &1.90mA\\
%       34    &26.8 &2.24mA\\
%       36    &26.1 &3.08mA\\
%       38    &25.6 &3.80mA\\
%       40    &25.1 &4.10  \\
%       42    &24.8 &5.34mA\\
%       45    &24.3 &6.41mA \\
%       48.7  &23.9 &7.62mA  
%     \end{tabular}
%   \end{center} 
% \end{table}

\begin{minipage}[t]{0.5\linewidth}
  \begin{table}[H] %R = 1000ohm
    \begin{center}
      \begin{tabular}{c|c|c}
        $V_{CC}$ &$V_{A_1A_2}$ &$I$  \\
        \hline
        5     &5    &0     \\
        10    &10   &1    uA\\
        15    &15   &1.5  uA\\
        20    &20   &2    uA\\
        22    &22   &2.2  uA\\
        24    &24   &2.4  uA\\
        26    &26   &2.6  uA\\
        28    &28   &2.7  uA\\
        30    &30   &3    uA\\
        32    &32   &3.2  uA\\
        33.4  &23.2 &10.5 mA\\
        34    &23.2 &11.1 mA\\
        35    &23.0 &12.3 mA\\
        35.5  &22.9 &12.8 mA\\
        36    &22.8 &13.5 mA\\
        37    &22.7 &14.7 mA\\
        38    &22.6 &15.8 mA\\
        40    &22.5 &18.1 mA\\
        42    &22.2 &20.7 mA\\
        45    &22   &23.7 mA \\
        50    &22.1 &29.1 mA \\
        55    &22.3 &34.6 mA \\
        60    &22.5 &40.2 mA \\
        63    &22.6 &43.5 mA \\
      \end{tabular}
    \end{center} 
  \end{table}
\end{minipage}
\begin{minipage}[t]{0.5\linewidth}
  \begin{table}[H] %R = 1000ohm
    \begin{center}
      \begin{tabular}{c|c|c}
        $V_{CC}$ &$V_{A_1A_2}$ &$I$  \\
        \hline
        -5    &-5    &-0     \\
        -10    &-10   &-1    uA\\
        -15    &-15   &-1.5  uA\\
        -20    &-20   &-2    uA\\
        -22    &-22   &-2.2  uA\\
        -24    &-24   &-2.4  uA\\
        -26    &-26   &-2.6  uA\\
        -28    &-28   &-2.7  uA\\
        -30    &-30   &-3    uA\\
        -32    &-32   &-3.2  uA\\
        -33    &-23   &-10.36 mA\\
        -34    &-22.7 &-11.53 mA\\
        -36    &-22.4 &-13.72 mA\\
        -38    &-22.2 &-16.07 mA\\
        -40    &-22.0 &-18.16 mA\\
        -42    &-21.9 &-20.60 mA\\
        -45    &-21.8 &-24.10 mA \\
        -50    &-21.7 &-29.50 mA \\
        -55    &-21.9 &-34.40 mA \\
        -60    &-22.1 &-39.70  mA \\
        -63    &-22.3 &-43.70  mA \\
      \end{tabular}
    \end{center} 
  \end{table}
\end{minipage}

% \begin{tikzpicture}
%   \begin{axis}[
%     xlabel = $V_{A_1A_2}$,
%     ylabel = $I$,
%     grid=major,
%     ]
%     \addplot[color=red,mark=*] coordinates {
%         (5   ,   0   )
%         (10  ,   1e-6 )
%         (15  ,   1.5e-6)
%         (20  ,   2e-6  )
%         (22  ,   2.2e-6)
%         (24  ,   2.4e-6)
%         (26  ,   2.6e-6)
%         (28  ,   2.7e-6)
%         (30  ,   3e-6  )
%         (32  ,   3.2e-6)
%         (27.2,   1.90e-3)
%         (26.8,   2.24e-3)
%         (26.1,   3.08e-3)
%         (25.6,   3.80e-3)
%         (25.1,   4.10e-3)
%         (24.8,   5.34e-3)
%         (24.3,   6.41e-3)
%         (23.9,   7.62e-3)
%       };  
%   \end{axis} 
% \end{tikzpicture}


\begin{tikzpicture}
  \begin{axis}[
    width=0.8\linewidth,
    xlabel = $V_{A_1A_2}$,
    ylabel = $I$,
    grid=major,
    ]
    \addplot[color=red,mark=*, mark size=1.2pt] coordinates {
        (-22.3 ,  -43.70e-3)
        (-22.1 ,  -39.70e-3)
        (-21.9 ,  -34.40e-3)
        (-21.7 ,  -29.50e-3)
        (-21.8 ,  -24.10e-3)
        (-21.9 ,  -20.60e-3)
        (-22.0 ,  -18.16e-3)
        (-22.2 ,  -16.07e-3)
        (-22.4 ,  -13.72e-3)
        (-22.7 ,  -11.53e-3)
        (-23   ,  -10.36e-3)
        (-32   ,  -3.2e-6  )
        (-30   ,  -3  e-6  )
        (-28   ,  -2.7e-6  )
        (-26   ,  -2.6e-6  )
        (-24   ,  -2.4e-6  )
        (-22   ,  -2.2e-6  )
        (-20   ,  -2e-6    )
        (-15   ,  -1.5e-6  )
        (-10   ,  -1e-6    )
        (-5    ,   0       )
        (5     ,   0       )
        (10    ,   1e-6    )
        (15    ,   1.5e-6  )
        (20    ,   2e-6    )
        (22    ,   2.2e-6  )
        (24    ,   2.4e-6  )
        (26    ,   2.6e-6  )
        (28    ,   2.7e-6  )
        (30    ,   3e-6    )
        (32    ,   3.2e-6  )
        (23.2  ,   10.5e-3 )
        (23.2  ,   11.1e-3 )
        (23.0  ,   12.3e-3 )
        (22.9  ,   12.8e-3 )
        (22.8  ,   13.5e-3 )
        (22.7  ,   14.7e-3 )
        (22.6  ,   15.8e-3 )
        (22.5  ,   18.1e-3 )
        (22.2  ,   20.7e-3 )
        (22    ,   23.7e-3 )
        (22.1  ,   29.1e-3 )
        (22.3  ,   34.6e-3 )
        (22.5  ,   40.2e-3 )
        (22.6  ,   43.5e-3 )
      };  
  \end{axis} 
\end{tikzpicture}
\section{Segunda Parte}
\subsection{Circuito}
\begin{center}
  \begin{circuitikz}[american]
    \draw (0,0) node[vcc](vcc){44V}
      to[R=$R_1$, a=3.3\unit{\kilo\ohm}] ++(0,-1.5)
      to[pR, n=pot, a=5\unit{\kilo\ohm}] ++(0,-4)
      to[short] ++(0,-0.5) node[ground](gnd1){};
    \draw (6,0) node[vcc](VCC){$V_{CC}$}
      to[R=$R_4$, a=5.8\unit{\kilo\ohm}] ++(0,-2.1)
      to[Ty*,n=tr,mirror, v^=$V_{AK}$,a=BT151-800R] (VCC |- pot.wiper)
      to[short,f^=$I_{AK}$] (VCC|-gnd1) node[ground]{};
    \draw (tr.gate) to[nos, mirror, invert] ++(-2,0) 
      to[R=$R_3$, a=1\unit{\kilo\ohm}] (pot.wiper);
  \end{circuitikz}
\end{center}
\subsection{Procedimiento}
\begin{enumerate}
  \item Armar el circuito.
  \item Colocar la VCC = 0V.
  \item Cerrar el interruptor Sw.
  \item Variar el potenciómetro de forma de relevar la tabla.
  \item Graficar los valores obtenidos y comparar la curva con la de otro
    componente ya estudiado.
  \item Abrir el interruptor Sw.
  \item Colocar un voltímetro en paralelo a la resistencia de carga (RL) y otro en paralelo al Anodo-Catodo del SCR.
  \item Variar la VCC desde 0V a 600V en pasos de 10V controlando permanentemente lo que sucede en los voltímetros.
  \item Finalizado el ensayo, ¿noto un cambio de comportamiento en el circuito?, ¿En qué valor de tensión?.
  \item Desconectando las alimentaciones de tensión, ¿puede analizar el valor ohmico de la 
    resistencia de carga (RL)?, ¿Qué sucedió?.
  \item Colocar la VG = 0V y cerrar el interruptor Sw.
  \item Colocar VCC = 100V.
  \item Subir lentamente el valor de VG hasta observar un cambio importante en la IAK (Disparo del SCR). 
    Tomar nota del valor de VG e IG que produjo ese disparo del SCR.
  \item Manteniendo el potenciómetro en la posición donde generó el disparo abrir el 
    interruptor Sw y analizar que sucede con la IAK.
  \item Manteniendo el interruptor Sw abierto bajar el valor de VCC en pasos de 10V 
    anotando el valor de IAK para cada caso. Los últimos 10V antes de llegar a cero deben disminuirse en pasos de 1V.
  \item Volver a subir paulatinamente la VCC hasta colocarla nuevamente en 100V
    analizando el comportamiento que tiene la IAK
  \item Cerrar el interruptor Sw corroborando que los valores de VG e IG sean los
mismos que se relevaron en el punto 13.
  \item Analizar que sucedió en IAK cuando se cerró el interruptor Sw.
  \item Abrir nuevamente el interruptor Sw y volver a observar la IAK.
  \item Desconectar las tensiones de alimentación sin mover la posición donde
quedo el potenciómetro e invertir el Ánodo y Cátodo del SCR.
  \item Volver a conectar las alimentaciones y cerrar el interruptor Sw. Analizar el
comportamiento en IAK.
\end{enumerate}
\subsection{Simulación}
\subsubsection{Comportamiento de la compuerta}
\begin{figure}[H]
  \begin{center}
    %\includegraphics[width=0.5\textwidth, clip, trim=0.5cm 10cm 0.5cm 8cm]{inc/sim/circuits/CircuitoSegundaParte.pdf}
  \end{center}
  \caption{Circuito simulado}
  \label{fig:CircuitoSegundaParte}
\end{figure}
\begin{center}
  \begin{tikzpicture}[
    ]
    \begin{axis}[
      width=\plotScale\linewidth,
      xlabel=$V_{G}$, 
      ylabel=$I_G$, 
      grid=both,
      scaled y ticks=base 10:3,
      thick
      ]
      \addplot[color=red,
        %skip coords between index={0}{0},
        ] table [x index=1, y index=2]{./inc/sim/SegundaParteA.txt};
    \end{axis}
  \end{tikzpicture}
\end{center}
\subsubsection{Disparó por tensión}
\begin{figure}[H]
  \begin{center}
    %\includegraphics[width=0.5\textwidth, clip, trim=0.5cm 10cm 0.5cm 8cm]{inc/sim/circuits/CircuitoSegundaParte.pdf}
  \end{center}
  \caption{Circuito simulado}
  \label{fig:CircuitoSegundaParte}
\end{figure}
\begin{center}
  \begin{tikzpicture}[
    ]
    \begin{axis}[
      width=\plotScale\linewidth,
      xlabel=$V_{AK}$, 
      ylabel=$I_{AK}$, 
      grid=both,
      scaled y ticks=base 10:3,
      thick
      ]
      \addplot[color=red,
        %skip coords between index={0}{0},
        ] table [x index=1, y index=3]{./inc/sim/SegundaParteB.txt};
    \end{axis}
  \end{tikzpicture}
\end{center}

\subsection{Experimental} %Alimentamos con 30V el div ressitivo
Colocando $V_{CC}=0V$ 
\par 
\begin{minipage}{0.3\linewidth}
  \centering
  \begin{table}[H]
    \begin{center}
      \begin{tabular}{c|c}
        $V_{G}$ &$I_{G}$[mA]   \\
        \hline
        0.1   &0.377 \\
        0.2   &0.722 \\
        0.3   &1.062 \\
        0.4   &1.40  \\
        0.5   &2.60  \\
        0.6   &3.08  \\
        0.7   &3.80  \\
        0.73  &4.40  \\
      \end{tabular}
    \end{center} 
  \end{table}
\end{minipage}
\begin{minipage}{0.7\linewidth}
  \centering
  \begin{tikzpicture}
    \begin{axis}[
      xlabel = $V_{G}$,
      ylabel = $I_G$,
      grid=major,
      ]
      \addplot[color=red,mark=*] coordinates {
          (0      ,   0)
          (0.05  ,0.023)
          (0.1   ,0.377)
          (0.2   ,0.722)
          (0.3   ,1.062)
          (0.4   ,1.40 )
          (0.5   ,2.50 )
          (0.6   ,3.08 )
          (0.7   ,3.80 )
          (0.73  ,4.40 )
          (0.76  ,4.87 )
        };  
    \end{axis} 
  \end{tikzpicture}
\end{minipage}
%V lado del pote = 44V
%Vcc = 208V
%Vgt = 1V 
%Igt = 3.38mA 
%Ih  =   
%Vt  =  1V 
%Ranodo = 5.8k
Posteriormente, se colocó $V_{CC}=210V$ y se disparó el SRC tomando las mediciones pertinentes. 
Luego del disparó, se interumpió $I_G$ y se procedió a disminuir $V_{CC}$ progresivamente hasta apagar el dispositivo.
\par
\begin{minipage}{0.3\linewidth}
\begin{table}[H]
  \begin{center}
    \begin{tabular}{c|c}
      $V_{CC}$ &$I_{AK}$[mA]   \\
      \hline
      195   &28   \\
      187   &26  \\
      180   &25.5  \\
      175   &24.6  \\
      170   &23.8  \\
      160   &22.7  \\
      149   &21  \\
      140   &19  \\
      110   &12.8  \\
      100   &11.7  \\
      90    &10.5  \\
      80    &9.5  \\
    \end{tabular}
  \end{center} 
\end{table}
\end{minipage}
\begin{minipage}{0.7\linewidth}
  \centering
  \begin{tikzpicture}
    \begin{axis}[
      xlabel = $V_{AK}$,
      ylabel = $I_{AK}$,
      grid=major,
      ]
      \addplot[color=red,mark=*] coordinates {
          (195 ,28  )  
          (187 ,26  )
          (180 ,25.5)
          (175 ,24.6)
          (170 ,23.8)
          (160 ,22.7)
          (149 ,21  )
          (140 ,19  )
          (110 ,12.8)
          (100 ,11.7)
          (90  ,10.5)
          (80  ,9.5 )
          (80  ,0.5 )
        };    
    \end{axis} 
  \end{tikzpicture}
\end{minipage}
\par
Posteriormente, luego del apagado del SCR, se volvió a conectar la corriente de compuerta,
manteniendo el potenciómetro en la posición previa. Analizando de este modo, que el SCR 
se encendió nuevamente.
\par
Por último, con el potenciómetro en la posición de disparo, se invertió la conexiones 
anodo-catodo del SCR. Así, se logró medir una corriente de fuga inversa de 12.5uA
\begin{align}
  I_{RR} = 12.5\unit{\micro\ampere} 
\end{align}
De esta manera, las mediciones tomadas en este último procedicimento resultan
\begin{table}[H]
  \begin{center}
    \begin{tabular}{l|c|c|c}
      \multirow{2}{*}{Parametro} &\multirow{2}{*}{Medido @$Vcc=210V$}  &\multicolumn{2}{c}{Datasheet}  \\
      \cline{3-4}
                &  &TYP. &MAX  \\
      \hline
      $V_{GT}$  &1V      &0.6V   &1.5V  \\
      $I_{GT}$  &3.38mA  &2mA   &15mA  \\
      $V_{T}$   &1V      &1.4V   &1.75V @$I_T=23A$  \\
      $I_{H}$   &9.50mA  &7mA   &20mA  \\
      $I_{R}$   &12.5\unit{\micro\ampere} & &
    \end{tabular}
  \end{center} 
\end{table}
Vale la pena aclarar que, a no ser que se aclare lo contrario, las conidicones para las cuales se especificaron los 
valores brindados por el fabricante (datasheet) son:
\begin{itemize}
  \item $T_j=25ºC$ 
  \item $V_D=12V$ 
  \item $I_{GT}=0.1A$ 
  \item $I_{T}=0.1A$ 
\end{itemize}

\section{Tercera Parte}
\subsection{Circuito}
Para este tercera parte, se trabajó con el mismo circuito que en la parte anterior.
\subsection{Procedimiento}
\begin{enumerate}
  \item Armar el circuito.
  \item Completar los valores de la tabla 1 fijando el valor de IG y variando el valor
    de VCC hasta observar el disparo del SCR. Relevar el valor de VCC, IAK y VAK
    en el cual se produce el disparo para cada caso 
  \item Completar los valores de la tabla 2 fijando el valor de IG y variando el valor
    de VCC. Relevar IAK y VAK para cada caso. Es posible que deba alterar
    algunos valores de resistencia dependiendo del SCR utilizado.
\end{enumerate}
\subsection{Experimental}
%44V del lado del pote
% \begin{table}[H]
%   \begin{center}
%     \begin{tabular}{c|c|c|c}
%       $I_{G}[\unit{\milli\ampere}]$ &$V_{CC}$ &$I_{AK}$ &$V_{AK}$ \\
%       \hline
%       4    &14  &5.70   &0.69 \\
%       4.5  &2   &5.20   &0.66 \\
%     \end{tabular}
%   \end{center} 
% \end{table}

\begin{minipage}{0.3\linewidth}
\begin{table}[H]
  \begin{center}
    \begin{tabular}{|l|c|c|c|c|}
      \hline
      $I_{G}\rightarrow$  &\multicolumn{2}{c|}{3.35 mA} &\multicolumn{2}{c|}{3.65 mA}  \\ 
      \hline
      \hline
      $V_{CC}$  &$V_{AK}$ &$I_{AK}$ &$V_{AK}$ &$I_{AK}$ \\ 
      \hline
      0     &0     &3.40    &0    &3.6    \\
      1     &0.60  &3.50    &0.6  &3.62   \\
      2     &1     &3.55    &0.73 &3.70   \\
      3     &1.60  &3.57    &1.09 &3.90   \\
      4     &2.33  &3.60    &1.1  &3.90   \\
      5     &2.46  &3.72    &1.33 &4      \\
      6     &2.50  &3.75    &1.35 &4.10   \\
      7     &2.81  &3.9     &1.25 &4.12   \\
      8     &0.67  &4.62    &0.75 &4.36   \\
      9     &0.69  &5.60    &0.79 &5.30   \\
      \hline
    \end{tabular}
  \end{center}
\end{table}
\end{minipage}
\begin{minipage}{0.7\linewidth}
\centering
  \begin{tikzpicture}
    \begin{axis}[
      xlabel = $V_{AK}$,
      ylabel = $I_{AK}$,
      scaled y ticks=base 10:3,
      grid=major,
      ]
      \addplot[color=red,mark=*] coordinates {
          (0    ,3.40 )  
          (0.60 ,3.50 )
          (1    ,3.55 )
          (1.60 ,3.57 )
          (2.33 ,3.60 )
          (2.46 ,3.72 )
          (2.60 ,3.75 )
          (2.81 ,3.9  )
          (0.67 ,4.62 )
          (0.69 ,5.62 )
        };    
      \addplot[color=blue,mark=*] coordinates {
          (0    ,3.6 )  
          (0.6  ,3.62)
          (0.73 ,3.70)
          (1.09 ,3.90)
          (1.1  ,3.90)
          (1.33 ,4   )
          (1.35 ,4.10)
          (1.25 ,4.12)
          (0.75 ,4.36)
          (0.78 ,5.36)
        };
    \end{axis} 
  \end{tikzpicture}
\end{minipage}


\section{Cuarta Parte}
\subsection{Circuito}
\begin{center}
  \begin{circuitikz}[american]
    \draw (0,0) node[vcc](vcc){20V}
      to[R=$R_1$, a=3.3\unit{\kilo\ohm}] ++(0,-1.5)
      to[pR, n=pot, a=5\unit{\kilo\ohm}] ++(0,-4)
      to[short] ++(0,-0.5) node[ground](gnd1){};
    \draw (6,0) node[vcc](VCC){$V_{CA}$}
      to[R=$R_4$, a=5.8\unit{\kilo\ohm}] ++(0,-2.1)
      to[Ty*,n=tr,mirror, v^=$V_{AK}$,a=BT151-800R] (VCC |- pot.wiper)
      to[R=$R5$,a=100\unit{\ohm},f^=$I_{AK}$] (VCC|-gnd1) node[ground]{};
    \draw (tr.gate) to[nos, mirror, invert] ++(-2,0) 
      to[R=$R_3$, a=4.7\unit{\kilo\ohm}] (pot.wiper);
  \end{circuitikz}
\end{center}
\subsection{Procedimiento}
\begin{enumerate}
  \item Armar el circuito.
  \item Variar el potenciómetro de forma que VG quede a potencial de cero volts
    cuando conectemos la fuente de alimentación.
  \item Colocar la tensión de alimentación alterna en 50 VCC
  \item Aumentar lentamente la VG observando permanentemente la IG e IA.
    Determinar el momento donde el dispositivo se dispara.
  \item Bajar el valor de VG a cero y observar lo que sucede con la IA.
  \item En función de lo estudiado en el teórico apagar el SCR.
  \item Subir ahora el valor de VCC a 100V y repetir los pasos 4, 5 y 6.
  \item Subir ahora el valor de VCC a 150V y repetir los pasos 4, 5 y 6.
  \item Completar la siguiente tabla de mediciones.
\end{enumerate}
\subsection{Simulación}
\subsection{Experimental}
\begin{table}[H]
  \begin{center}
    \begin{tabular}{|l|c|c|c|c|c|c|}
      \hline
      $V_{CC}\rightarrow$ &\multicolumn{2}{|c|}{$V_{CC}=50$} &\multicolumn{2}{c|}{$V_{CC}=100$} &\multicolumn{2}{c|}{$V_{CC}=150$} \\ 
      \hline
      \hline
      $V_{G}$ &$I_G$ &$I_A$ &$I_G$ &$I_A$  &$I_G$ &$I_A$ \\ 
      \hline
      0     &0     &0    &0    &0     &0     &0      \\
      0.25  &0.90  &0.90 &0.88 &0.88  &0.86  &0.86   \\
      0.50  &1.72  &1.72 &1.83 &1.83  &1.80  &1.80   \\
      0.75  &2.64  &2.64 &2.60 &2.60  &2.58  &2.58   \\
      1     &3.52  &3.52 &3.50 &3.50  &3.48  &3.48   \\
      1.40  &3.52  &6.56 &-    &-     &-     &-      \\
      1.50  &4.45  &7.22 &-    &-     &-     &-      \\
      1.55  &4.93  &7.56 &-    &-     &-     &-      \\
      1.70  &-     &-    &3.30 &9.36  &-     &-      \\
      1.70  &-     &-    &3.77 &9.86  &-     &-      \\
      1.80  &-     &-    &7.74 &10.47 &-     &-      \\
      1.90  &-     &-    &-    &-     &-     &-      \\
      2     &-     &-    &-    &-     &3.11  &12     \\
      2.10  &-     &-    &-    &-     &3.40  &12.60  \\
      2.20  &-     &-    &-    &-     &4.40  &13.39  \\
      \hline
    \end{tabular}
  \end{center}
\end{table}

\section{Quinta Parte}
\subsection{Circuito}
%potencimetro 5k
\begin{center}
  \begin{circuitikz}[american]
    \ctikzset{bipoles/border margin=1}
    \draw (0,0) node[vcc](vcc){$V_{AC}$} 
      to[R=$R_6$, a=5\unit{\kilo\ohm}]++(0,-3) coordinate(R6)
      to[Tr*, mirror, n=triac, l=BT136]++(0,-4)
      to[short, -*]++(0,-1) node[ground](gnd){};
    \draw (triac.G) to[biD*, a=DB3]++(-2,0) --++(-0.5,0) node[potentiometershape, anchor=wiper, rotate=-90, name=5\unit{\kilo\ohm}](pot){5\unit{\kilo\ohm}};
    \draw (pot) to[short]++(0,2) coordinate(A) to[short,-*](R6|-A);
    \draw (pot) to[short]++(0,-2) coordinate(B) to[short,-*](gnd|-B);
  \end{circuitikz}
\end{center}
\subsection{Procedimiento}
\begin{enumerate}
  \item Armar el circuito.
  \item Variar el potenciómetro de forma que VG quede a potencial de cero volts
    cuando conectemos la fuente de alimentación.
  \item  Colocar la tensión de alimentación alterna en 50 VCC
  \item Aumentar lentamente la VG observando permanentemente la IG e IA. Determinar el momento donde el dispositivo se dispara.
  \item Bajar el valor de VG a cero y observar lo que sucede con la IA.
  \item En función de lo estudiado en el teórico apagar el TRIAC.
  \item Subir ahora el valor de VCC a 100V y repetir los pasos 4, 5 y 6.
  \item Subir ahora el valor de VCC a 150V y repetir los pasos 4, 5 y 6.
  \item Completar la siguiente tabla de mediciones.
\end{enumerate}
\subsection{Simulación}
%Usamos 5k
\subsection{Experimental}
\begin{table}[H]
  \begin{center}
    \begin{tabular}{|c|c|c|c|c|c|c|c|c|}
      \hline
      $V_{CC}\rightarrow $&\multicolumn{2}{|c|}{$V_{AC}=50$} &\multicolumn{2}{c|}{$V_{AC}=75$} &\multicolumn{2}{c|}{$V_{AC}=90$} &\multicolumn{2}{c|}{$V_{AC}=105$}\\ 
      \hline
      \hline
      $V_{1}$ &$I_G$ &$I_A$  &$I_G$ &$I_A$  &$I_G$ &$I_K$  &$I_G$ &$I_K$\\ 
      \hline
      0    &0      &0        &176u    &10.2m   &100u   &15.6m &79.8u    &18.8m    \\
      5    &5.10u  &5.10u    &548u    &7.08m   &268u   &12.0m &139.7u   &16.09m   \\
      10   &9.80u  &9.80u    &224u    &3.85m   &-      &-     &300.2u   &13.44m   \\
      15   &15.10u &15.10u   &14.9u   &14.9u   &473u   &8.35m &-        &-        \\
      20   &20.60u &20.60u   &20.3u   &20.3u   &20.2u  &20.2u &-        &-        \\
      23   &22.80u &22.80u   &23.1u   &23.1u   &23.1u  &23.1u &440u     &9.32m    \\
      30   &-      &-        &30.4u   &30.4u   &30.4u  &30.4u &30.5u    &30.5u    \\
      34   &-      &-        &34.6u   &34.6u   &34.6.u &34.6u &34.9u    &34.9u    \\
      40   &-      &-        &-       &-       &40.6u  &40.6u &40.6u    &40.6u    \\
            \hline
    \end{tabular}
  \end{center}
\end{table}
De esta manera los momentos en los que el triac se dispara son:
\begin{enumerate}
  \item Para @$V_{CC}=50V$
    \begin{center}
      No se disparó
    \end{center}
  \item Para @$V_{CC}=75V$ 
    \begin{align}
      V_{GT} &= 10V - V_{T_{\text{Diac}}}  \\
      I_{GT} &= 188u \\
      I_{A}  &= 0.25mA 
    \end{align}
  \item Para @$V_{CC}=90V$ 
    \begin{align}
      V_{GT} &= 16.8V - V_{T_{\text{Diac}}}  \\
      I_{GT} &= 430u \\
      I_{A}  &=  8.07mA 
    \end{align}
  \item Para @$V_{CC}=105V$ 
    \begin{align}
      V_{GT} &= 24.3V - V_{T_{\text{Diac}}}  \\
      I_{GT} &= 223u \\
      I_{A}  &= 8.50mA 
    \end{align}
\end{enumerate}
\begin{center}
  \begin{tikzpicture}
    \begin{axis}[
      xlabel = $I_{G}$,
      ylabel = $I_{A}$,
      grid=major,
      scaled y ticks=base 10:3,
      legend pos=outer north east
      ]
      \addplot[color=red, mark=*] coordinates {
          (0       ,0       )  
          (5.10e-6 ,5.10e-6 )
          (9.80e-6 ,9.80e-6 )
          (15.10e-6,15.10e-6)
          (20.60e-6,20.60e-6)
          (22.80e-6,22.80e-6)
        };    
      \addplot[color=blue, mark={square*}] coordinates {
          (176e-6  ,  10e-3)  
          (548e-6  ,7.08e-3)
          (224e-6  ,3.85e-3)
          (14.9e-6 ,14.9e-6)
          (20.3e-6 ,20.3e-6)
          (23.1e-6 ,23.1e-6)
          (30.4e-6 ,30.4e-6)
          (34.6e-6 ,34.6e-6)
        };
      \addplot[color=brown, mark={triangle*}] coordinates {
          (100e-6   ,15.6e-3)  
          (268e-6   ,12.0e-3)
          (473e-6   ,8.35e-3)
          (20.2e-6  ,20.2e-6)
          (23.1e-6  ,23.1e-6)
          (30.4e-6  ,30.4e-6)
          (34.6e-6  ,34.6e-6)
          (40.6e-6  ,40.6e-6)
        };         
     \addplot[color=violet, mark=x] coordinates {
          (79.8e-6   ,18.8e-3  )  
          (139.7e-6  ,16.09e-3 )
          (300.2e-6  ,13.44e-3 )
          (440e-6    , 9.32e-3 )
          (30.5e-6   , 30.5e-6 )
          (34.9e-6   , 34.9e-6 )
          (40.6e-6   , 40.6e-6 )
        };              
      \legend{
          $V_{CC}=50V$,
          $V_{CC}=75V$,
          $V_{CC}=90V$,
          $V_{CC}=105V$,
        }
    \end{axis}          
  \end{tikzpicture}
\end{center}

\section{Séptima Parte}
\subsection{Circuito}
\begin{center}
  \begin{circuitikz}[american]
    \ctikzset{bipoles/border margin=1}
    \draw (0,0) node[vcc](vcc){$V_{AC}$} 
      to[R=$R_6$, a=5\unit{\kilo\ohm}]++(0,-3) coordinate(R6)
      to[Tr*, mirror, n=triac, l=BT136]++(0,-4)
      to[R,a=100\unit{\ohm}]++(0,-2) node[ground](gnd){};
    \draw (triac.G) to[biD*, a=DB3]++(-2,0) to[R,a=220\unit{\ohm}]++(-2,0) --++(-0.5,0) node[potentiometershape, anchor=wiper, rotate=-90, name=5\unit{\kilo\ohm}](pot){5\unit{\kilo\ohm}};
    \draw (pot) to[R,a=2.2\unit{\kilo\ohm}]++(0,3) coordinate(A) to[short,-*](R6|-A);
    \draw (pot) to[C,a=0.1\unit{\micro\farad}]++(0,-3) coordinate(B) to[short,-*](gnd|-B);
  \end{circuitikz}
\end{center}
\subsection{Procedimiento}
\begin{itemize}
  \item  Procedimiento
  \item Armar el circuito.
  \item Llevar el potenciómetro hacia un extremo y medir la forma de onda en RL y
  luego entre ambos ánodos del TRIAC.
  \item Colocar el potenciómetro en su posición media y medir la forma de onda en
  RL y luego entre ambos ánodos del TRIAC.
  \item Colocar el potenciómetro hacia el otro extremo y medir la forma de onda en
  RL y luego entre ambos ánodos del TRIAC.
\end{itemize}
\subsection{Simulación}
\begin{center}
  \begin{tikzpicture}
    \begin{groupplot}[
      group style={ group size=1 by 3, x descriptions at=edge bottom, vertical sep=3pt},
      xlabel=$t$,
      width=\plotScale\linewidth,
      height=0.5\linewidth,
      grid=both,
      legend pos=outer north east
      ]
      \nextgroupplot[ylabel=$V_{DIAC}$]
      \addplot[thick, color=red, mark={triangle},mark repeat=80,mark size=1.5pt,
        %skip coords between index={0}{291},
        skip coords between index={291}{1196}
        ] table [x index=0, y index=1]{./inc/sim/SeptimaParte.txt};
      \addplot[thick,color=blue, mark={square},mark repeat=80,mark size=1.5pt,
        skip coords between index={0}{291},
        skip coords between index={670}{1196}
        ] table [x index=0, y index=1]{./inc/sim/SeptimaParte.txt};
      \addplot[thick,color=brown, mark=x,mark repeat=80,mark size=1.5pt,
        skip coords between index={0}{670}
        ] table [x index=0, y index=1]{./inc/sim/SeptimaParte.txt};
     \legend{
     $K\rightarrow 1$,
     $K\rightarrow 0.5$,
     $K\rightarrow 0$
     };
      \nextgroupplot[ylabel=$V_{A_1}$]
      \addplot[thick,color=red, mark={triangle},mark repeat=80,mark size=1.5pt,
        %skip coords between index={0}{291},
        skip coords between index={291}{1196},
        ] table [x index=0, y expr=\thisrowno{2}-\thisrowno{3}]{./inc/sim/SeptimaParte.txt};
      \addplot[thick,color=blue, mark={square},mark repeat=80,mark size=1.5pt,
        skip coords between index={0}{291},
        skip coords between index={670}{1196},
        ] table [x index=0, y expr=\thisrowno{2}-\thisrowno{3}]{./inc/sim/SeptimaParte.txt};
      \addplot[thick,color=brown, mark=x,mark repeat=80,mark size=1.5pt,
        skip coords between index={0}{670},
        ] table [x index=0, y expr=\thisrowno{2}-\thisrowno{3}]{./inc/sim/SeptimaParte.txt};
      \nextgroupplot[ylabel=$V_{L}$]
      \addplot[thick,color=red, mark={triangle},mark repeat=80,mark size=1.5pt,
        %skip coords between index={0}{291},
        skip coords between index={291}{1196},
        ] table [x index=0, y index=3]{./inc/sim/SeptimaParte.txt};
      \addplot[thick,color=blue, mark={square},mark repeat=80,mark size=1.5pt,
        skip coords between index={0}{291},
        skip coords between index={670}{1196},
        ] table [x index=0, y index=3]{./inc/sim/SeptimaParte.txt};
      \addplot[thick,color=brown, mark=x,mark repeat=80,mark size=1.5pt,
        skip coords between index={0}{670},
        ] table [x index=0, y index=3]{./inc/sim/SeptimaParte.txt};
    \end{groupplot}
  \end{tikzpicture}
\end{center}

\section{Octava Parte}
\subsection{Circuito}
\begin{center}
  \begin{circuitikz}[american]
    \ctikzset{bipoles/border margin=1}
    \draw (0,0) node[vcc](vcc){$V_{AC}$} 
      to[L, a=2H]++(0,-2) to[R=$R_6$, a=5\unit{\kilo\ohm}]++(0,-3) coordinate(R6)
      to[Tr*, mirror, n=triac, l=BT136]++(0,-4)
      to[R,a=100\unit{\ohm}]++(0,-2) node[ground](gnd){};
    \draw (triac.G) to[biD*, a=DB3]++(-2,0) to[R,a=220\unit{\ohm}]++(-2,0) --++(-0.5,0) node[potentiometershape, anchor=wiper, rotate=-90, name=5\unit{\kilo\ohm}](pot){5\unit{\kilo\ohm}};
    \draw (pot) to[R,a=2.2\unit{\kilo\ohm}]++(0,3) coordinate(A) to[short,-*](R6|-A);
    \draw (pot) to[C,a=0.1\unit{\micro\farad}]++(0,-3) coordinate(B) to[short,-*](gnd|-B);
  \end{circuitikz}
\end{center}
\subsection{Procedimiento}
\begin{itemize}
  \item Armar el circuito.
  \item  Cerrar el interruptor de forma que la parte inductiva de la carga quede
     cortocircuitada.
  \item Colocar el potenciómetro en su posición media entre ambos ánodos del
    TRIAC. Observar y graficar con detenimiento el momento donde la IH
    produce el apagado del TRIAC.
  \item Abrir el interruptor de forma que la parte inductiva de la carga quede
    conectada.
  \item Colocar el potenciómetro en su posición media entre ambos ánodos del
    TRIAC. Observar y graficar con detenimiento el momento donde la IH
    produce el apagado del TRIAC
\end{itemize}
\subsection{Simulación}
\begin{figure}[H]
  \begin{center}
    %\includegraphics[width=0.95\textwidth]{}
  \end{center}
  \caption{Circuito simulado}\label{fig:}
\end{figure}
\subsubsection{Sin inductor}
\begin{center}
  \begin{tikzpicture}
    \begin{groupplot}[
      group style={ group size=1 by 3, x descriptions at=edge bottom, vertical sep=3pt},
      xlabel=$t$,
      width=\plotScale\linewidth,
      height=0.5\linewidth,
      grid=both,
      legend pos=outer north east
      ]
      \nextgroupplot[ylabel=$V_{DIAC}$]
      \addplot[thick, color=red, mark={triangle},mark repeat=80,mark size=1.5pt,
        %skip coords between index={0}{291},
        skip coords between index={291}{1196}
        ] table [x index=0, y index=1]{./inc/sim/OctaveParte.txt};
      \addplot[thick,color=blue, mark={square},mark repeat=80,mark size=1.5pt,
        skip coords between index={0}{291},
        skip coords between index={670}{1196}
        ] table [x index=0, y index=1]{./inc/sim/OctaveParte.txt};
      \addplot[thick,color=brown, mark=x,mark repeat=80,mark size=1.5pt,
        skip coords between index={0}{670}
        ] table [x index=0, y index=1]{./inc/sim/OctaveParte.txt};
     \legend{
     $K\rightarrow 1$,
     $K\rightarrow 0.5$,
     $K\rightarrow 0$
     };
      \nextgroupplot[ylabel=$V_{A_1}$]
      \addplot[thick,color=red, mark={triangle},mark repeat=80,mark size=1.5pt,
        %skip coords between index={0}{291},
        skip coords between index={291}{1196},
        ] table [x index=0, y expr=\thisrowno{2}-\thisrowno{3}]{./inc/sim/OctaveParte.txt};
      \addplot[thick,color=blue, mark={square},mark repeat=80,mark size=1.5pt,
        skip coords between index={0}{291},
        skip coords between index={670}{1196},
        ] table [x index=0, y expr=\thisrowno{2}-\thisrowno{3}]{./inc/sim/OctaveParte.txt};
      \addplot[thick,color=brown, mark=x,mark repeat=80,mark size=1.5pt,
        skip coords between index={0}{670},
        ] table [x index=0, y expr=\thisrowno{2}-\thisrowno{3}]{./inc/sim/OctaveParte.txt};
      \nextgroupplot[ylabel=$V_{L}$]
      \addplot[thick,color=red, mark={triangle},mark repeat=80,mark size=1.5pt,
        %skip coords between index={0}{291},
        skip coords between index={291}{1196},
        ] table [x index=0, y index=3]{./inc/sim/OctaveParte.txt};
      \addplot[thick,color=blue, mark={square},mark repeat=80,mark size=1.5pt,
        skip coords between index={0}{291},
        skip coords between index={670}{1196},
        ] table [x index=0, y index=3]{./inc/sim/OctaveParte.txt};
      \addplot[thick,color=brown, mark=x,mark repeat=80,mark size=1.5pt,
        skip coords between index={0}{670},
        ] table [x index=0, y index=3]{./inc/sim/OctaveParte.txt};
    \end{groupplot}
  \end{tikzpicture}
\end{center}
\subsubsection{Con inductor}
\begin{center}
  \begin{tikzpicture}
    \begin{groupplot}[
      group style={ group size=1 by 3, x descriptions at=edge bottom, vertical sep=3pt},
      xlabel=$t$,
      width=\plotScale\linewidth,
      height=0.5\linewidth,
      grid=both,
      legend pos=outer north east
      ]
      \nextgroupplot[ylabel=$V_{DIAC}$]
      \addplot[thick, color=red, mark={triangle},mark repeat=80,mark size=1.5pt,
        %skip coords between index={0}{291},
        skip coords between index={269}{1073}
        ] table [x index=0, y index=1]{./inc/sim/OctaveParteInductor.txt};
      \addplot[thick,color=blue, mark={square},mark repeat=80,mark size=1.5pt,
        skip coords between index={0}{269},
        skip coords between index={617}{1073}
        ] table [x index=0, y index=1]{./inc/sim/OctaveParteInductor.txt};
      \addplot[thick,color=brown, mark=x,mark repeat=80,mark size=1.5pt,
        skip coords between index={0}{617}
        ] table [x index=0, y index=1]{./inc/sim/OctaveParteInductor.txt};
     \legend{
     $K\rightarrow 1$,
     $K\rightarrow 0.5$,
     $K\rightarrow 0$
     };
      \nextgroupplot[ylabel=$V_{A_1}$]
      \addplot[thick,color=red, mark={triangle},mark repeat=80,mark size=1.5pt,
        %skip coords between index={0}{291},
        skip coords between index={269}{1073}
        ] table [x index=0, y expr=\thisrowno{2}-\thisrowno{3}]{./inc/sim/OctaveParteInductor.txt};
      \addplot[thick,color=blue, mark={square},mark repeat=80,mark size=1.5pt,
        skip coords between index={0}{269},
        skip coords between index={617}{1073}
        ] table [x index=0, y expr=\thisrowno{2}-\thisrowno{3}]{./inc/sim/OctaveParteInductor.txt};
      \addplot[thick,color=brown, mark=x,mark repeat=80,mark size=1.5pt,
        skip coords between index={0}{617}
        ] table [x index=0, y expr=\thisrowno{2}-\thisrowno{3}]{./inc/sim/OctaveParteInductor.txt};
      \nextgroupplot[ylabel=$V_{L}$]
      \addplot[thick,color=red, mark={triangle},mark repeat=80,mark size=1.5pt,
        %skip coords between index={0}{291},
        skip coords between index={269}{1073}
        ] table [x index=0, y index=3]{./inc/sim/OctaveParteInductor.txt};
      \addplot[thick,color=blue, mark={square},mark repeat=80,mark size=1.5pt,
        skip coords between index={0}{269},
        skip coords between index={617}{1073}
        ] table [x index=0, y index=3]{./inc/sim/OctaveParteInductor.txt};
      \addplot[thick,color=brown, mark=x,mark repeat=80,mark size=1.5pt,
        skip coords between index={0}{617}
        ] table [x index=0, y index=3]{./inc/sim/OctaveParteInductor.txt};
    \end{groupplot}
  \end{tikzpicture}
\end{center}



\end{document}
